% ═══════════════════════════════════════════════════════════════════
% Group Chat Application - Lab 4 Report
% Secure Persistent Group Chat: Encryption, Signing & Database Layer
% Builds upon the Lab 3 WebSocket Chat (report.tex)
% ═══════════════════════════════════════════════════════════════════

\documentclass[12pt, a4paper]{article}

\usepackage[utf8]{inputenc}
\usepackage[T1]{fontenc}
\usepackage{lmodern}
\usepackage[margin=1in]{geometry}
\usepackage{graphicx}
\usepackage{hyperref}
\usepackage{xcolor}
\usepackage{listings}
\usepackage{booktabs}
\usepackage{enumitem}
\usepackage{float}
\usepackage{caption}
\usepackage{fancyhdr}
\usepackage{titlesec}
\usepackage{tocloft}
\usepackage{amsmath}
\usepackage{parskip}
\usepackage{tabularx}
\usepackage{longtable}
\usepackage{multirow}

\definecolor{codebg}{HTML}{F5F5F5}
\definecolor{codeframe}{HTML}{CCCCCC}
\definecolor{keyword}{HTML}{0077AA}
\definecolor{comment}{HTML}{999999}
\definecolor{string}{HTML}{669900}
\definecolor{linkblue}{HTML}{0066CC}

\hypersetup{
    colorlinks=true,
    linkcolor=linkblue,
    urlcolor=linkblue,
    citecolor=linkblue,
    pdfauthor={Team Members},
    pdftitle={PixelChat Lab 4 - Secure Persistent Group Chat},
    pdfsubject={Computer System Design Tutorial - Lab 4},
}

\lstdefinestyle{codestyle}{
    backgroundcolor=\color{codebg},
    frame=single,
    rulecolor=\color{codeframe},
    basicstyle=\ttfamily\footnotesize,
    keywordstyle=\color{keyword}\bfseries,
    commentstyle=\color{comment}\itshape,
    stringstyle=\color{string},
    breaklines=true,
    breakatwhitespace=false,
    numbers=left,
    numberstyle=\tiny\color{comment},
    numbersep=8pt,
    showstringspaces=false,
    tabsize=4,
    captionpos=b,
}
\lstset{style=codestyle}

\pagestyle{fancy}
\fancyhf{}
\fancyhead[L]{\small PixelChat -- Lab 4: Secure Persistent Group Chat}
\fancyhead[R]{\small \thepage}
\fancyfoot[C]{\small Computer System Design Tutorial Report}
\renewcommand{\headrulewidth}{0.4pt}
\renewcommand{\footrulewidth}{0.4pt}

\titleformat{\section}{\Large\bfseries}{\thesection.}{0.5em}{}
\titleformat{\subsection}{\large\bfseries}{\thesubsection.}{0.5em}{}
\titleformat{\subsubsection}{\normalsize\bfseries}{\thesubsubsection.}{0.5em}{}

\begin{document}

\begin{titlepage}
    \centering
    \vspace*{1cm}

    {\LARGE\bfseries Report [Lab 4] }\\[0.5cm]
    {\Huge\bfseries Group Chat Application using web sockets}\\[0.8cm]
    {\large\itshape End-to-End Encryption, Digital Signatures, Database Persistence \& Multi-Room Architecture}\\[1.2cm]

    \rule{\textwidth}{0.5pt}\\[0.5cm]

    {\large Subject: Computer System Design}\\[0.3cm]
    {\large Course Code: \texttt{CSL559}}\\[0.3cm]
    {\large Semester: \texttt{7}}\\[0.5cm]

    \rule{\textwidth}{0.5pt}\\[1.2cm]

    {\large\bfseries Team Members}\\[0.5cm]

    \begin{tabular}{lll}
        \toprule
        \textbf{Name} & \textbf{Roll No.}\\
        \midrule
        Sneha Nagmoti& 12342090 \\
        Snehal Suhane & 12342100 \\
        Sannidhi Rithika & 12341900  \\
        Harshita Patidar & 12340920  \\
        \bottomrule
    \end{tabular}
   

\begin{center}
    \url{https://github.com/snehanagmoti/group-chat-app}
\end{center}

    \vfill

  

\end{titlepage}

\newpage
\tableofcontents
\newpage

%=================================================================
\section{Introduction}
%=================================================================

This report describes the design and implementation of the \textbf{Lab 4 extension} to the PixelChat Group Chat Application originally built in Lab 3. The Lab 3 report (\texttt{report.tex}) covers the foundational WebSocket-based real-time chat with avatars, gamification, and file attachments. This second phase adds enterprise-grade security, persistence, and advanced features as required by the Lab 4 specification.

\subsection{What Was Built in Lab 3}

The Lab 3 codebase established:
\begin{itemize}[leftmargin=*]
    \item Real-time WebSocket group chat with broadcast messaging to all connected users.
    \item 8-bit retro ``PixelChat'' UI (Press Start 2P font, pixel-art theme).
    \item Avatar picker (12 themed characters), emoji picker, typing indicators.
    \item File attachments (image preview, video/audio players, file cards).
    \item Message delivery receipts (sent / partial / delivered\_all).
    \item Auto-reconnect with exponential backoff on connection loss.
    \item Gamification: XP bar, rank progression (Newbie to Legend), streaks, session timer.
    \item Mario-themed sound effects for join/leave/message events.
\end{itemize}

\subsection{What Lab 4 Adds}

Lab 4 extends the above with the following major pillars:

\begin{itemize}[leftmargin=*]
    \item \textbf{User Authentication} -- bcrypt-hashed password accounts, one-time 256-bit session tokens.
    \item \textbf{AES-256-GCM End-to-End Encryption} -- all message content encrypted in the browser; the server never sees plaintext.
    \item \textbf{ECDSA-P256 Digital Signatures} -- every message signed by its author; server verifies each signature.
    \item \textbf{SQLite Database Persistence} -- messages, users, public keys, and rooms durably stored across restarts.
    \item \textbf{HMAC-SHA256 Tamper Detection} -- each database row integrity-protected; tampering detected and flagged.
    \item \textbf{Multi-Room Architecture} -- unlimited independent rooms with unique 6-character codes.
    \item \textbf{TLS/SSL Transport} -- all HTTP and WebSocket traffic encrypted with a self-signed TLS certificate.
    \item \textbf{Advanced Messaging} -- threaded replies, whisper/private messages, message unsend, 5-minute edit window.
    \item \textbf{Voice Message E2EE} -- microphone recordings encrypted with AES-GCM before transmission.
    \item \textbf{E2EE File Attachments} -- files encrypted before upload; decrypted by recipients.
    \item \textbf{Persistent XP} -- experience points stored in SQLite, survive server restarts.
    \item \textbf{Multi-Tab Session Sharing} -- same browser window tabs share login; new windows require fresh login.
    \item \textbf{Port Forwarding Support} -- frontend/backend port separation via environment variables.
\end{itemize}

\subsection{Lab 4 Objectives}

\begin{itemize}[leftmargin=*]
    \item Implement persistent user accounts with secure password hashing (bcrypt).
    \item Encrypt all chat messages end-to-end using AES-256-GCM via the browser SubtleCrypto API.
    \item Sign every message with an ECDSA-P256 browser-generated keypair; verify on the server.
    \item Persist encrypted, signed messages to SQLite with HMAC-SHA256 tamper detection.
    \item Implement multi-room support with creator-managed rooms.
    \item Demonstrate the server stores only ciphertext and can detect tampered database entries.
    \item Deploy with TLS: frontend on port 3000, backend on port 5000 (frontend accessible via port forwarding at port 3269)
\end{itemize}

%=================================================================
\section{Updated Technology Stack}
%=================================================================

\begin{table}[H]
    \centering
    \caption{Lab 4 Technology Stack}
    \begin{tabularx}{\textwidth}{l l X}
        \toprule
        \textbf{Layer} & \textbf{Technology} & \textbf{Purpose} \\
        \midrule
        Backend Server   & Python 3 + FastAPI + Uvicorn & Async WebSocket server, REST API, TLS-enabled \\
        Frontend Server  & Python 3 + FastAPI + Uvicorn & TLS static file server \\
        Database         & SQLite 3 (built-in)          & Persistent storage: users, rooms, encrypted messages, keys \\
        Password Hashing & \texttt{bcrypt}              & Slow-hash; resistant to brute-force \\
        Encryption       & AES-256-GCM (SubtleCrypto)   & Client-side E2EE of all message content \\
        Digital Signing  & ECDSA-P256 / SHA-256         & Per-user signing; browser keypair, server verification \\
        Tamper Detection & HMAC-SHA256                  & Row-level DB integrity \\
        TLS Transport    & OpenSSL self-signed cert     & Encrypted HTTPS/WSS transport \\
        SQL Layer        & Python \texttt{sqlite3}      & Direct parameterized SQL; no external ORM \\
        Env Config       & \texttt{python-dotenv}       & Secrets loaded from \texttt{.env}; never hardcoded \\
        \bottomrule
    \end{tabularx}
\end{table}

\subsection{Python Dependencies (requirements.txt)}

\begin{table}[H]
    \centering
    \caption{Python Dependencies}
    \begin{tabularx}{\textwidth}{l X}
        \toprule
        \textbf{Package} & \textbf{Purpose} \\
        \midrule
        \texttt{fastapi}           & Web framework: WebSocket, REST, file upload support \\
        \texttt{uvicorn[standard]} & ASGI server; TLS via ssl\_keyfile / ssl\_certfile \\
        \texttt{python-dotenv}     & Load \texttt{.env} secrets at startup \\
        \texttt{python-multipart}  & Multipart parsing for file uploads \\
        \texttt{bcrypt}            & Industry-standard password hashing (Blowfish) \\
        \texttt{cryptography}      & ECDSA-P256 public key import and signature verification \\
        \bottomrule
    \end{tabularx}
\end{table}


%=================================================================
\section{System Architecture}
%=================================================================

\subsection{Architecture Overview}



Lab 4 extends Lab 3with five additional security/persistence layers:

\begin{enumerate}[leftmargin=*]
    \item \textbf{TLS Transport Layer} -- Frontend (port 3000) and backend (port 5000) both run with a self-signed TLS certificate. All communication is HTTPS/WSS.
    \item \textbf{Authentication Layer} -- Browser POSTs to \texttt{/register} or \texttt{/login}; server returns a one-time 256-bit cryptographic token consumed on WebSocket connect.
    \item \textbf{Client-Side Encryption \& Signing} -- Browser SubtleCrypto generates an ECDSA-P256 keypair per session. Each message is AES-GCM encrypted then ECDSA signed before transmission. Only ciphertext, IV, signature, and JWK public key travel the network.
    \item \textbf{Server-Side Verification} -- Server reconstructs the public key from JWK and calls \texttt{verify\_ecdsa\_signature()}. Invalid signatures flagged in logs and UI.
    \item \textbf{SQLite Persistence \& HMAC} -- Each message persisted with HMAC-SHA256(ciphertext||iv). On history load, HMAC recomputed; mismatch $\Rightarrow$ \texttt{tampered=True} $\Rightarrow$ warning in UI.
\end{enumerate}

\subsection{Port Configuration and Forwarding}

\begin{table}[H]
    \centering
    \caption{Port Mapping for Lab Deployment}
    \begin{tabularx}{\textwidth}{l c c X}
        \toprule
        \textbf{Service} & \textbf{Env Variable} & \textbf{Port} & \textbf{Notes} \\
        \midrule
        Frontend Web Server   & \texttt{FRONTEND\_PORT}    & 3000 & Browser opens https://host:3000 \\
        Backend API/WS Server & \texttt{PORT}              & 5000 & FastAPI/Uvicorn binds here \\
        Browser to Backend    & \texttt{BACKEND\_PORT}     & 5269 & Port forwarding: 5269 $\to$ 5000 \\
        \bottomrule
    \end{tabularx}
\end{table}

The \texttt{/config.js} endpoint (served by the frontend server) injects \texttt{window.BACKEND\_PORT = 5269} into the browser, directing all API and WebSocket requests to port 5269 (the forwarded port), while the backend actually listens on 5000 (via the \texttt{PORT} environment variable).


%=================================================================
\section{User Authentication}
%=================================================================

\subsection{Overview}

Lab 4 introduces persistent user accounts backed by SQLite. Users must register or log in before accessing any room. Authentication is \textbf{stateless on the server} after login: a one-time token is issued and consumed exactly once on WebSocket connect.

\subsection{Registration (POST /register)}

\begin{enumerate}[leftmargin=*]
    \item Client sends \texttt{\{username, password, avatar\}}.
    \item Server validates: username 1--20 alphanumeric chars; password $\geq$ 6 chars.
    \item Password hashed with \texttt{bcrypt}: \texttt{bcrypt.hashpw(password, bcrypt.gensalt())}.
    \item Row inserted into \texttt{users} table (UNIQUE constraint on username).
    \item A one-time session token (\texttt{secrets.token\_hex(32)}, 256 bits) returned.
\end{enumerate}
\begin{figure}[H]
    \centering
    \includegraphics[width=0.85\textwidth]{3.png}
    \caption{Sign in}
    \label{fig:chat}
\end{figure}
\subsection{Login (POST /login)}

\begin{enumerate}[leftmargin=*]
    \item Client sends \texttt{\{username, password\}}.
    \item Server retrieves user row (case-insensitive lookup).
    \item \texttt{bcrypt.checkpw()} verifies the password.
    \item On success: token issued, stored in in-memory \texttt{active\_sessions}; green log printed.
    \item On failure: HTTP 401 returned; red error log printed.
\end{enumerate}


\begin{figure}[H]
    \centering
    \includegraphics[width=0.85\textwidth]{2.png}
    \caption{Login}
    \label{fig:chat}
\end{figure}

\subsection{WebSocket Token Handshake}

The one-time token is sent in the first WebSocket message (\texttt{type: "join"}). The server calls \texttt{active\_sessions.pop(token, None)} -- consuming it permanently. If absent or invalid, the socket closes immediately (code 1008).

\begin{lstlisting}[language=Python, caption={Token Validation in WebSocket Endpoint (server.py)}]
token   = data.get("token", "").strip()
pub_key = data.get("public_key")   # ECDSA JWK
room_id = data.get("room_id", "").strip().upper()

session = active_sessions.pop(token, None)  # one-time use
if not session:
    await websocket.send_json({"type": "error",
                                "message": "Invalid or expired token."})
    await websocket.close(code=1008)
    return
\end{lstlisting}



\subsection{Authentication Console Logs}

The backend prints colored terminal logs for all authentication events:

\begin{figure}[H]
    \centering
    \includegraphics[width=0.85\textwidth]{13.jpeg}
    \caption{Authentication Console Logs}
    \label{fig:chat}
\end{figure}



%=================================================================
\section{End-to-End Encryption (AES-256-GCM)}
%=================================================================

\subsection{Overview}

All message content -- text, voice memos, and file attachments -- is encrypted \textbf{in the browser} using AES-256-GCM before being sent. The server receives, stores, and forwards only ciphertext. Decryption happens in the recipient's browser. The server never has access to plaintext.

\subsection{Key Distribution}

A single 256-bit AES-GCM group key is shared by all room participants. It is stored in \texttt{.env} and served to clients via \texttt{GET /group-key} after login.

\begin{lstlisting}[language=Python, caption={Group Key REST Endpoint (server.py)}]
@app.get("/group-key")
async def get_group_key():
    return {"key": AES_GROUP_KEY_HEX}  # loaded from .env at startup
\end{lstlisting}

The browser imports the key via SubtleCrypto:

\begin{lstlisting}[language=JavaScript, caption={AES Key Import in Browser (app.js)}]
const keyBytes = hexToBytes(keyHex);   // 32 bytes
AES_KEY = await crypto.subtle.importKey(
    "raw", keyBytes.buffer,
    { name: "AES-GCM" },
    false,                             // non-extractable
    ["encrypt", "decrypt"]
);
\end{lstlisting}

\subsection{Encryption Process}

\begin{enumerate}[leftmargin=*]
    \item A fresh 12-byte IV is generated with \texttt{crypto.getRandomValues(new Uint8Array(12))}.
    \item Message text (UTF-8 encoded) encrypted with AES-256-GCM + group key + IV.
    \item Ciphertext and IV base64-encoded and sent as JSON fields.
\end{enumerate}

\begin{lstlisting}[language=JavaScript, caption={Message Encryption (app.js)}]
async function encryptMessage(plaintext) {
    const iv = crypto.getRandomValues(new Uint8Array(12));
    const buf = await crypto.subtle.encrypt(
        { name: "AES-GCM", iv },
        AES_KEY,
        new TextEncoder().encode(plaintext)
    );
    return { ciphertext: bufToBase64(buf), iv: bufToBase64(iv) };
}
\end{lstlisting}

\subsection{Decryption Process}

On receiving a message (live or from history):
\begin{enumerate}[leftmargin=*]
    \item Decode base64 ciphertext and IV.
    \item Call \texttt{crypto.subtle.decrypt()} with group AES key and message IV.
    \item On failure: show \texttt{[DECRYPTION FAILED]} warning on the message bubble.
\end{enumerate}

\subsection{Voice Message Encryption}

\begin{enumerate}[leftmargin=*]
    \item Audio blob from \texttt{MediaRecorder} converted to base64 Data URL.
    \item Data URL string AES-GCM encrypted (fresh IV).
    \item Sent as \texttt{type: "voice"} message with ciphertext/IV fields.
    \item Recipients decrypt and render an inline HTML5 audio player.
\end{enumerate}

\subsection{File / Media Attachment Encryption}

\begin{enumerate}[leftmargin=*]
    \item File read as ArrayBuffer in the browser.
    \item Raw bytes AES-GCM encrypted.
    \item Encrypted Blob uploaded via \texttt{POST /upload}.
    \item Server stores only the encrypted file. Recipients download and decrypt.
    \item A compact \texttt{lock and tick} badge on the attachment card shows E2EE status.
\end{enumerate}

\begin{figure}[H]
    \centering
    \includegraphics[width=0.85\textwidth]{16.jpeg}
    \caption{File encryption}
    \label{fig:chat}
\end{figure}

%=================================================================
\section{Digital Signatures (ECDSA-P256)}
%=================================================================

\subsection{Overview}

Every message is cryptographically signed by its author using an \textbf{ECDSA-P256} keypair. The signature proves authenticity: only the private-key holder could produce it. The server and all recipients can verify using the corresponding public key.

\subsection{Keypair Generation}

At login, the browser generates a fresh ECDSA-P256 keypair:

\begin{lstlisting}[language=JavaScript, caption={ECDSA Keypair Generation (app.js)}]
const keyPair = await crypto.subtle.generateKey(
    { name: "ECDSA", namedCurve: "P-256" },
    true,                  // public key exportable
    ["sign", "verify"]
);
ECDSA_PRIVATE_KEY = keyPair.privateKey;
const jwk = await crypto.subtle.exportKey("jwk", keyPair.publicKey);
// jwk sent to server in the WebSocket join message
\end{lstlisting}

The JWK public key is registered in the \texttt{user\_keys} SQLite table and included with every message.

\subsection{Message Signing}

The client signs the concatenation of \texttt{ciphertext + iv} (UTF-8 bytes):

\begin{lstlisting}[language=JavaScript, caption={Per-Message Signing (app.js)}]
async function signMessage(ciphertext, iv) {
    const material = new TextEncoder().encode(ciphertext + iv);
    const sigBuf = await crypto.subtle.sign(
        { name: "ECDSA", hash: { name: "SHA-256" } },
        ECDSA_PRIVATE_KEY,
        material
    );
    return bufToBase64(sigBuf);  // IEEE P1363: 64 bytes for P-256
}
\end{lstlisting}

\subsection{Server-Side Verification}

The server verifies every incoming message using the Python \texttt{cryptography} library:

\begin{lstlisting}[language=Python, caption={ECDSA Verification (server.py)}]
def verify_ecdsa_signature(plaintext_bytes, sig_b64, jwk):
    pub_key = _jwk_to_public_key(jwk)   # rebuild from JWK
    sig_bytes = base64.urlsafe_b64decode(sig_b64 + "=" * (-len(sig_b64) % 4))
    # Convert IEEE P1363 (r||s, 64 bytes) -> DER for the Python library
    r = int.from_bytes(sig_bytes[:32], "big")
    s = int.from_bytes(sig_bytes[32:], "big")
    sig_der = encode_dss_signature(r, s)
    pub_key.verify(sig_der, plaintext_bytes, ECDSA(hashes.SHA256()))
    return True   # raises InvalidSignature on failure

# In the WebSocket message loop:
signed_material = (ciphertext + iv).encode("utf-8")
sig_valid = verify_ecdsa_signature(signed_material, signature, sender_key)
if sig_valid:
    print(f"\033[92m[AUTH VERIFIED] ECDSA-P256 VERIFIED for '@{username}'\033[0m")
else:
    print(f"\033[91m[SECURITY ALERT] INVALID SIGNATURE from '@{username}'\033[0m")
\end{lstlisting}

\begin{figure}[H]
    \centering
    \includegraphics[width=0.85\textwidth]{auth.jpeg}
    
    
 \caption{ECDSA-P256 Signature Verification -- Valid  and Invalid }
    \label{fig:chat}
\end{figure}
\begin{figure}[H]
    \centering
    \includegraphics[width=0.85\textwidth]{unauth.jpeg}
    
    
    \caption{ECDSA-P256 Signature Verification -- Valid  and Invalid }
    \label{fig:chat}
\end{figure}

\subsection{Signature Result Persisted in Database}

The \texttt{sig\_valid} boolean is stored in the \texttt{messages} table. When a user joins and receives history, pre-computed \texttt{sig\_valid} flags are included. Messages with \texttt{sig\_valid = 0} render a warning indicator.

%=================================================================
\section{Database Persistence (SQLite)}
%=================================================================

\subsection{Overview}

All application data is persisted in \texttt{server/chat.db}, initialised at startup via \texttt{db.init\_db()}. The database layer is encapsulated in \texttt{server/db.py}.

\subsection{Schema: Four Tables}

\begin{table}[H]
    \centering
    \small
    \caption{SQLite Database Schema}
    \begin{tabularx}{\textwidth}{l X X}
        \toprule
        \textbf{Table} & \textbf{Key Columns} & \textbf{Purpose} \\
        \midrule
        \texttt{messages}   & msg\_id, room\_id, username, ciphertext, iv, signature, public\_key, hmac\_digest, sig\_valid & Encrypted, signed, HMAC-protected messages \\
        \texttt{users}      & username (UNIQUE), password\_hash, avatar, xp & Accounts with bcrypt hashes and XP \\
        \texttt{user\_keys} & username (PK), public\_key (JWK) & ECDSA public key registry \\
        \texttt{rooms}      & id (6-char), name, created\_by, is\_public, avatar & Chat room directory \\
        \bottomrule
    \end{tabularx}
\end{table}



\subsection{Database Migrations}

The \texttt{\_apply\_migrations()} function uses \texttt{PRAGMA table\_info} to detect missing columns and applies \texttt{ALTER TABLE ADD COLUMN} statements, allowing schema evolution without database deletion.


    \begin{figure}[H]
    \centering
    \includegraphics[width=0.85\textwidth]{pers.jpeg}
    
    
   
    \caption{Message Persistence}
    
\end{figure}
  \begin{figure}[H]
    \centering
    \includegraphics[width=0.85\textwidth]{pers2.jpeg}
    
    
   
    \caption{Message Persistence -- Database}
    
\end{figure}

%=================================================================
\section{HMAC-SHA256 Tamper Detection}
%=================================================================

\subsection{Overview}

Every message row is protected by an \textbf{HMAC-SHA256} digest. This allows the server to detect if a row has been modified directly in the database file without knowledge of the server-side \texttt{HMAC\_SECRET}.

\subsection{HMAC at Save Time}

\begin{lstlisting}[language=Python, caption={HMAC Computation -- db.py}]
def _compute_hmac(ciphertext: str, iv: str) -> str:
    secret  = bytes.fromhex(os.environ["HMAC_SECRET"])  # 32 bytes
    payload = (ciphertext + iv).encode("utf-8")
    return hmac.new(secret, payload, hashlib.sha256).hexdigest()

# In save_message():
hmac_digest = _compute_hmac(ciphertext, iv)
# Stored in hmac_digest column alongside ciphertext
\end{lstlisting}

\subsection{HMAC at Load Time}

\begin{lstlisting}[language=Python, caption={HMAC Verification on History Load -- db.py}]
def _verify_hmac(ciphertext: str, iv: str, stored_digest: str) -> bool:
    expected = _compute_hmac(ciphertext, iv)
    return hmac.compare_digest(expected, stored_digest)  # constant-time

# In get_history():
tampered = not _verify_hmac(row["ciphertext"], row["iv"], row["hmac_digest"])
# tampered=True included in JSON sent to client; UI shows warning
\end{lstlisting}

\subsection{Security Properties of HMAC Layer}

\begin{itemize}[leftmargin=*]
    \item The \texttt{HMAC\_SECRET} (256 bits) is stored only in \texttt{.env}; never transmitted.
    \item \texttt{hmac.compare\_digest()} prevents timing-based oracle attacks.
    \item Tampered messages are flagged visibly, not silently suppressed -- preserving auditability.
    \item HMAC provides integrity independently from AES-GCM authentication tags.
\end{itemize}

  \begin{figure}[H]
    \centering
    \includegraphics[width=0.85\textwidth]{beforetamp.jpeg}
    
    
   
    \caption{Normal Message Flow}
    
\end{figure}
 \begin{figure}[H]
    \centering
    \includegraphics[width=0.85\textwidth]{tamper.jpeg}
    
    
   
    \caption{Tampering the message}
    
\end{figure}
  \begin{figure}[H]
    \centering
    \includegraphics[width=0.85\textwidth]{17.jpeg}
    
    
   
    \caption{After tampering}
    
\end{figure}

%=================================================================
\section{Multi-Room Architecture}
%=================================================================

\subsection{Room Management}

The application supports unlimited independent chat rooms. Each room is identified by a unique 6-character alphanumeric code (e.g., \texttt{XKJ3P9}), generated cryptographically with \texttt{secrets.choice()}.

\begin{table}[H]
    \centering
    \caption{Room Management REST Endpoints}
    \begin{tabularx}{\textwidth}{l l X}
        \toprule
        \textbf{Method \& Path} & \textbf{Auth} & \textbf{Description} \\
        \midrule
        GET /rooms              & None & List all public rooms with live online count \\
        POST /rooms             & Body & Create room; returns code and XP awarded \\
        GET /rooms/\{id\}       & None & Check if room exists; return metadata \\
        DELETE /rooms/\{id\}    & Creator & Permanently delete room and all messages \\
        DELETE /rooms/\{id\}/history & Creator & Clear room message history \\
        \bottomrule
    \end{tabularx}
\end{table}

\subsection{Room Lobby Screen}

After login, users see the \textbf{room lobby}: a grid of public room cards showing name, avatar emoji, creator, and live online count. Users can browse, join by code, or create a new room.

 \begin{figure}[H]
    \centering
    \includegraphics[width=0.85\textwidth]{6.png}
    
    
   
   
    

    \caption{Room Lobby -- Browsing and Joining Chat Rooms}
    \label{fig:lobby}
\end{figure}

\subsection{Creator Badges and Privileges}

The user who clicks \textbf{CREATE ROOM} is recorded as \texttt{created\_by} in the \texttt{rooms} table. System-seeded rooms use \texttt{created\_by = "system"}; the first real user to join is promoted to creator via \texttt{db.update\_room\_creator\_if\_system()}.

Creator privileges:
\begin{itemize}[leftmargin=*]
    \item \textbf{CREATOR} badge displayed next to their name in the room.
    \item Exclusive ability to clear room message history.
    \item Exclusive ability to permanently delete the room.
\end{itemize}

 \begin{figure}[H]
    \centering
    \includegraphics[width=0.85\textwidth]{5.png}
    
    
   
   
    \caption{Room Creator Badge and Exclusive Room Management Privileges}
    \label{fig:creatorbadge}
\end{figure}

\subsection{Multi-Tab Multi-Room Support}

Users can have multiple browser tabs open in different rooms under the same login. \texttt{ConnectionManager.get\_room\_users()} deduplicates by username, so a user in two tabs appears once in the online list.

%=================================================================
\section{Advanced Messaging Features}
%=================================================================

\subsection{Threaded Replies}

Reply to a specific message by clicking the reply button on any bubble. The \texttt{reply\_to} field (parent \texttt{msg\_id}) is stored in the database. Recipients see a quoted preview of the parent message.

\subsection{Whisper / Private Messages}

Send a private whisper to a specific room member. The server routes the message only to that user (not broadcast). Stored with \texttt{target\_user} set; filtered from other users' history loads.

 \begin{figure}[H]
    \centering
    \includegraphics[width=0.85\textwidth]{8.png}
    
    
   
  
    \caption{Whisper (Private Message)}
    \label{fig:whisper}
\end{figure}

 \begin{figure}[H]
    \centering
    \includegraphics[width=0.85\textwidth]{reply.png}
    
    
   
  
    \caption{Threaded Reply)}
    \label{fig:whisper}
\end{figure}

\subsection{Message Unsend (Soft Delete)}

Users can unsend their own messages. The server performs a \textbf{soft delete}: \texttt{is\_deleted = 1}, ciphertext/IV/signature cleared. All room members receive \texttt{message\_deleted} event. Bubble shows ``[message deleted]''.

\subsection{Message Editing (5-Minute Window)}

Users can edit their own messages within 5 minutes of sending. The edit re-encrypts and re-signs the new content. The database row is updated atomically (new ciphertext, IV, signature, HMAC digest; \texttt{is\_edited = 1}). The UI shows an \textbf{(edited)} label. Edits after 5 minutes are rejected server-side.

 \begin{figure}[H]
    \centering
    \includegraphics[width=0.85\textwidth]{edit1.png}
    \caption{Original Message}
    \label{fig:whisper}
\end{figure}
 \begin{figure}[H]
    \centering
    \includegraphics[width=0.85\textwidth]{edit3.png}
    \caption{The edit(within 5 minute window)}
    \label{fig:whisper}
\end{figure}
 \begin{figure}[H]
    \centering
    \includegraphics[width=0.85\textwidth]{edit4.png}
    \caption{Edited Message}
    \label{fig:whisper}
\end{figure}


\subsection{Voice Messages}

\begin{enumerate}[leftmargin=*]
    \item User clicks mic button; browser records via \texttt{MediaRecorder} API.
    \item Recorded Blob converted to base64 Data URL.
    \item Data URL \textbf{AES-GCM encrypted} (fresh IV, same group key).
    \item Encrypted voice memo sent as \texttt{type: "voice"} message (ciphertext + IV + ECDSA signature).
    \item Recipients decrypt and render an inline HTML5 audio player.
\end{enumerate}

 \begin{figure}[H]
    \centering
    \includegraphics[width=0.85\textwidth]{19.jpeg}
 

    \caption{Encrypted Voice Message with Inline Audio Player}
    \label{fig:voice}
\end{figure}

%=================================================================
\section{Persistent XP System (Server-Backed)}
%=================================================================

XP moved from client-side in Lab 3 to server-side SQLite in Lab 4. XP is atomically incremented server-side and pushed to the client via \texttt{xp\_update} WebSocket events.

\begin{table}[H]
    \centering
    \caption{Server-Side XP Events (Lab 4)}
    \begin{tabularx}{\textwidth}{l c X}
        \toprule
        \textbf{Event} & \textbf{XP} & \textbf{Notes} \\
        \midrule
        Send a message       & +10 & Per message sent \\
        Create a room        & +20 & On room creation \\
        Someone joins        & +3  & Awarded to all existing room members \\
        Receive a message    & +2  & Awarded to all room members \\
        Streak bonus         & +25 & Every 10th consecutive message sent \\
        Heartbeat (per min.) & +5  & Client sends heartbeat every 60 seconds \\
        \bottomrule
    \end{tabularx}
\end{table}

\begin{figure}[H]
    \centering
    \includegraphics[width=0.85\textwidth]{10.png}
    \caption{Feature index}
    \label{fig:chat}
\end{figure}

%=================================================================
\section{TLS/SSL Transport Security}
%=================================================================

A self-signed TLS certificate (\texttt{cert.pem}) and private key (\texttt{key.pem}) generated via \texttt{generate\_certs.py}. Both frontend (port 3000, \texttt{FRONTEND\_PORT}) and backend (port 5000, \texttt{PORT}) start Uvicorn with \texttt{ssl\_keyfile} and \texttt{ssl\_certfile}, enabling HTTPS and WSS.

\begin{lstlisting}[language=Python, caption={Uvicorn TLS Configuration (server.py)}]
uvicorn.run(
    app,
    host="0.0.0.0",
    port=port,                                          # 5000
    ssl_keyfile=os.path.join(BASE_DIR, "key.pem"),
    ssl_certfile=os.path.join(BASE_DIR, "cert.pem")
)
\end{lstlisting}

Because the certificate is self-signed, browsers warn on first access. An in-app \textbf{CONNECTION BLOCKED} overlay guides users to accept the certificate at the backend health URL. The app resumes automatically once accepted.

 \begin{figure}[H]
    \centering
    \includegraphics[width=0.85\textwidth]{1.png}
   
    \caption{TLS Self-Signed Certificate Warning and In-App Authorization Guide}
    \label{fig:tls}
\end{figure}
 \begin{figure}[H]
    \centering
    \includegraphics[width=0.85\textwidth]{9.png}
   
    \caption{After authorizing}
    \label{fig:tls}
\end{figure}

%=================================================================
\section{Updated WebSocket Protocol}
%=================================================================

\subsection{New Client to Server Message Types}

\begin{table}[H]
    \centering
    \small % Optional: Slightly reduces font size for cleaner spacing
    \caption{Lab 4 -- New Client-to-Server Messages}
    \begin{tabularx}{\textwidth}{l X X}
        \toprule
        \textbf{Type} & \textbf{Key Fields} & \textbf{Description} \\
        \midrule
        join             & token, room\_id, public\_key & Authenticate and join a room \\
        message          & ciphertext, iv, signature, public\_key, client\_msg\_id, reply\_to, target\_user, attachment & Encrypted and signed chat message \\
        delete\_message  & msg\_id & Soft-delete own message \\
        edit\_message    & msg\_id, ciphertext, iv, signature, public\_key & Re-encrypt and re-sign own message \\
        clear\_room\_history & --- & Creator clears all room messages \\
        delete\_room     & --- & Creator permanently deletes room \\
        heartbeat        & --- & Time-in-room XP tick (every 60s) \\
        typing           & --- & Typing indicator (throttled) \\
        \bottomrule
    \end{tabularx}
\end{table}

\subsection{New Server to Client Message Types}

\begin{table}[H]
    \centering
    \small
    \caption{Lab 4 -- New Server-to-Client Messages}
    \begin{tabularx}{\textwidth}{l X X}
        \toprule
        \textbf{Type} & \textbf{Key Fields} & \textbf{Description} \\
        \midrule
        message              & ciphertext, iv, signature, sig\_valid, tampered & Encrypted broadcast; client decrypts \\
        history              & messages[] with tampered, sig\_valid per entry & Full room history from DB \\
        xp\_update           & xp, gained, reason & Server-computed XP push \\
        message\_deleted     & msg\_id, username & Tombstone for unsent messages \\
        message\_edited      & msg\_id, ciphertext, iv, sig\_valid & Updated ciphertext broadcast \\
        room\_deleted        & room\_id & Room deleted by creator \\
        room\_history\_cleared & room\_id & Room history cleared by creator \\
        \bottomrule
    \end{tabularx}
\end{table}

\subsection{Full Encrypted Message Flow}

\begin{lstlisting}[language={}]
1. Client (Alice) -> Server:
   { "type": "message",
     "ciphertext": "HqA3Bp...==",   <- AES-GCM encrypted text
     "iv":         "dXrK9A...==",   <- 12-byte IV (base64)
     "signature":  "MEUCIQDj...",   <- ECDSA-P256/SHA-256 (base64)
     "public_key": { "kty":"EC", "crv":"P-256", "x":"...", "y":"..." },
     "client_msg_id": "msg-17238...-abc" }

2. Server: verifies ECDSA signature -> sig_valid = True
           computes HMAC(ciphertext || iv) -> saves to SQLite

3. Server -> ALL Clients (broadcast):
   { "type": "message", "ciphertext": "HqA3Bp...==",
     "iv": "dXrK9A...==", "sig_valid": true, "tampered": false,
     "username": "Alice", "timestamp": "21:30:05" }

4. Each client's browser: crypto.subtle.decrypt() -> "Hello everyone!"
   Plaintext only ever visible inside the browser.

5. Server -> Alice (receipt):
   { "type": "receipt", "msg_id": "msg-17238...-abc",
     "status": "delivered_all" }
\end{lstlisting}

%=================================================================
\section{Testing (Lab 4)}
%=================================================================

\begin{longtable}{p{0.5cm} p{5.2cm} p{4.5cm} p{2.5cm}}
    \toprule
    \textbf{\#} & \textbf{Test Case} & \textbf{Expected Result} & \textbf{Status} \\
    \midrule
    \endhead
    1  & Register with valid credentials & Account created; token returned & Pass \\
    \midrule
    2  & Register with duplicate username & HTTP 409 error & Pass \\
    \midrule
    3  & Login with correct password & bcrypt verified; token issued; green log & Pass \\
    \midrule
    4  & Login with wrong password & HTTP 401; red error log & Pass \\
    \midrule
    5  & Send a text message & AES-GCM encrypted; ECDSA signed; DB saved & Pass \\
    \midrule
    6  & Recipient decrypts message & Plaintext rendered; badge shown & Pass \\
    \midrule
    7  & Server restart; rejoin room & Full message history loaded from SQLite & Pass \\
    \midrule
    8  & Modify ciphertext in DB directly & Tamper warning shown in UI & Pass \\
    \midrule
    9  & Modify signature in DB directly & SIG INVALID warning; SECURITY ALERT log & Pass \\
    \midrule
    10 & Create a new chat room & Room in lobby; unique 6-char code; +20 XP & Pass \\
    \midrule
    11 & Join room by code & History loaded; user appears in sidebar & Pass \\
    \midrule
    12 & Creator deletes room & All members notified; room removed & Pass \\
    \midrule
    13 & Non-creator tries to delete room & HTTP 403 error returned & Pass \\
    \midrule
    14 & Unsend own message & Tombstone broadcast; [message deleted] shown & Pass \\
    \midrule
    15 & Edit own message within 5 min & Re-encrypted; (edited) label shown & Pass \\
    \midrule
    16 & Edit after 5-minute window & Server rejects; error message returned & Pass \\
    \midrule
    17 & Whisper to a specific user & Only target user receives message & Pass \\
    \midrule
    18 & Record and send voice memo & Encrypted; audio player rendered & Pass \\
    \midrule
    19 & Upload and send file attachment & Encrypted upload; recipient decrypts & Pass \\
    \midrule
    20 & Open second tab (same window) & Auto-logged in with same session & Pass \\
    \midrule
    21 & Open new browser window & Login screen shown; no session shared & Pass \\
    \midrule
    22 & XP persists after server restart & XP total restored from DB at login & Pass \\
    \midrule
    23 & 4 users chat across lab machines & All encrypted messages delivered in real-time & Pass \\
    \midrule
    24 & TLS cert accepted; WSS established & Secure WebSocket active; padlock in browser & Pass \\
    \bottomrule
\end{longtable}

 \begin{figure}[H]
    \centering
    \includegraphics[width=0.85\textwidth]{log1.jpeg}
    \caption{Backend Terminal -- Full Authentication and Security Logging During Test Session}
    \label{fig:backlogs}
\end{figure}
 \begin{figure}[H]
    \centering
    \includegraphics[width=0.85\textwidth]{log2.jpeg}
    \caption{Backend Terminal -- Full Authentication and Security Logging During Test Session}
    \label{fig:backlogs}
\end{figure}
 \begin{figure}[H]
    \centering
    \includegraphics[width=0.85\textwidth]{log3.jpeg}
    \caption{Backend Terminal -- Full Authentication and Security Logging During Test Session}
    \label{fig:backlogs}
\end{figure}


%=================================================================
\section{Deployment Instructions (Lab 4)}
%=================================================================

\subsection{Prerequisites}

Python 3.10+; \texttt{pip}; OpenSSL (for cert generation); lab machines on same network with port forwarding (\texttt{BACKEND\_PORT}=5269 $\to$ \texttt{PORT}=5000); modern browser with SubtleCrypto API.

\subsection{Installation}

\begin{lstlisting}[language=bash, caption={Installation Steps}]
git clone https://github.com/snehanagmoti/group-chat-app && cd group-chat-app
cp .env.example .env            # Edit AES_GROUP_KEY, HMAC_SECRET, ports
python3 generate_certs.py       # Creates cert.pem and key.pem
pip install -r server/requirements.txt
\end{lstlisting}

\subsection{Starting the Servers}

\begin{lstlisting}[language=bash, caption={Starting Both Servers}]
# Terminal 1 -- Backend (PORT=5000)
python3 server/server.py
# Secure Group Chat Server (Multi-Room)
# [DB] Initialised -- server/chat.db

# Terminal 2 -- Frontend (FRONTEND_PORT=3000)
python3 client/serve.py
\end{lstlisting}

\subsection{Lab Deployment (4+ Machines)}

\begin{enumerate}[leftmargin=*]
    \item \textbf{Server machine:} Run both servers.
    \item \textbf{All clients:} Open \texttt{https://10.1.75.51:3269} in browser (\texttt{FRONTEND\_PORT}=3000).
    \item \textbf{Accept TLS cert:} Visit \texttt{https://10.1.75.51:5269/health} (\texttt{BACKEND\_PORT}=5269), accept warning, return.
    \item \textbf{Register/Login:} Each user creates an account.
    \item \textbf{Join or Create Room:} Select from lobby or enter a room code.
    \item \textbf{Chat Securely:} All messages AES-GCM encrypted and ECDSA signed end-to-end.
\end{enumerate}

%=================================================================
\section{Conclusion}
%=================================================================

This Lab 4 extension successfully transforms the Lab 3 WebSocket chat into a production-quality \textbf{secure, persistent group chat application}. Key achievements:

\begin{itemize}[leftmargin=*]
    \item[\checkmark] \textbf{AES-256-GCM E2EE} -- server never processes plaintext.
    \item[\checkmark] \textbf{ECDSA-P256 signatures} -- every message authenticated; server-verified and DB-persisted.
    \item[\checkmark] \textbf{HMAC-SHA256 tamper detection} -- row-level DB integrity; flagged in UI.
    \item[\checkmark] \textbf{SQLite persistence} -- users, rooms, messages, XP stored durably.
    \item[\checkmark] \textbf{bcrypt authentication} -- secure password storage and login.
    \item[\checkmark] \textbf{Multi-room architecture} -- unlimited rooms with creator controls.
    \item[\checkmark] \textbf{TLS transport} -- all traffic encrypted with HTTPS/WSS.
    \item[\checkmark] \textbf{Advanced messaging} -- whispers, replies, edit, unsend.
    \item[\checkmark] \textbf{Voice and file E2EE} -- media encrypted before transmission.
    \item[\checkmark] \textbf{Persistent XP} -- gamification state survives restarts.
\end{itemize}

The implementation demonstrates practical applied cryptography, secure protocol design, database persistence patterns, and WebSocket architecture -- all running in a standard browser without plugins.

%=================================================================

%=================================================================



%=================================================================

\section{Team Contributions}
%=================================================================

This section documents the individual contributions of each team member to the Lab 4 extension of PixelChat, covering the period after the Lab 3 baseline was completed. All four contributors participated across implementation, integration, and debugging. Contributions are grouped by feature area rather than individual commits.


%------------------------------------------------------------------
\subsection{Harshita Patidar}
\textit{GitHub: \texttt{harshitap1305}}

\textbf{Primary Areas:} Initial Persistence Integration, Encryption, User Authentication, Login/Registration UI, Documentation.

\begin{itemize}[leftmargin=*]
    \item \textbf{Persistence \& AES-GCM Encryption Integration:} Architected the secure database layer by integrating SQLite persistence with client-side AES-256-GCM encryption. Ensured a zero-knowledge storage model where only ciphertext and initialization vectors (IVs) are persisted. Engineered the secure distribution of the shared AES group key via a dedicated \texttt{/group-key} endpoint, allowing seamless decryption of chat history upon room join using the Web Crypto API (\texttt{SubtleCrypto}).

    \item \textbf{Tamper Detection \& Security Badges:} Designed and deployed an HMAC-SHA256 based cryptographic tamper detection system to ensure data integrity at rest. The system automatically computes and verifies message digests during database operations to intercept unauthorized modifications. Complemented this backend security with a frontend visual indicator system, displaying dynamic security badges (e.g., Verified, Invalid Signature, Tampered) to expose the cryptographic integrity of every message to the end user.

    \item \textbf{Secure User Authentication System:} Engineered a complete, end-to-end user authentication pipeline. Designed the SQLite \texttt{users} schema to track credentials, avatars, and experience points (XP). Developed the \texttt{POST /register} and \texttt{POST /login} REST endpoints utilizing \texttt{bcrypt} password hashing (cost factor 12) for secure credential storage, alongside a robust 256-bit cryptographic session token issuance system. Enforced rigorous server-side validation for username uniqueness and password complexity.

    \item \textbf{Responsive Login \& Registration UI:} Developed the interactive frontend login and registration experience. Implemented comprehensive client-side error handling for edge cases (e.g., duplicate usernames, invalid credentials, network timeouts) and built a dynamic avatar picker grid. Architected the layout to be fully responsive, identifying and resolving a critical CSS overflow bug to ensure seamless accessibility across small viewports and high browser zoom levels.

    \item \textbf{Multi-Tab Concurrency Fixes:} Resolved critical concurrency bugs arising from users maintaining multiple open browser tabs simultaneously. Engineered solutions to prevent duplicate username rendering in the active user list, fixed whisper message routing failures for targets with multiple active connections, and eliminated spurious join/leave broadcast events when closing redundant tabs, significantly improving system stability.

    \item \textbf{Deployment Configuration \& Documentation:} Standardized the environment deployment by correcting IP binding and port configurations for lab environments. Authored comprehensive inline docstrings for the authentication endpoints and database modules, and contributed extensive setup and deployment instructions to the project \texttt{README.md}.
\end{itemize}

%------------------------------------------------------------------
\subsection{Rithika Sannidhi}
\textit{GitHub: \texttt{rithika1503}}

\textbf{Primary Areas:} Room Management, Message Lifecycle, E2EE for Voice and Files, Security Logging, Port Decoupling, Session Management.

\begin{itemize}[leftmargin=*]
    \item \textbf{Room Management (Creator Controls):} Implemented the full room lifecycle on both server and client: \texttt{delete\_room} and \texttt{clear\_room\_history} WebSocket handlers, corresponding \texttt{db.delete\_room()} and \texttt{db.clear\_room\_history\_by\_creator()} database functions, and UI controls visible only to the room creator. Extended the room listing API to include the creator's avatar by joining the \texttt{users} table. Also disabled auto-cleanup timers to ensure messages are retained indefinitely in SQLite.

    \item \textbf{Message Lifecycle (Edit, Delete, History):} Implemented the message editing feature with a 5-minute server-side edit window enforced via \texttt{created\_at\_ts}. Edits re-encrypt and re-sign the new content, atomically replacing the HMAC digest in the database. Added the \texttt{(edited)} label in the UI. Removed the 50-message history cap, making the full SQLite history available on join. Also fixed UI state management (button resets, stale timers on tab switch).

    \item \textbf{E2EE for Voice Messages and File Attachments:} Integrated AES-256-GCM encryption into both the voice memo pipeline (MediaRecorder output encrypted before transmission; recipients receive a decrypted HTML5 audio player) and the file/media attachment pipeline (files read as ArrayBuffer, encrypted before upload, decrypted client-side on receipt; E2EE badge on attachment cards).

    \item \textbf{Security Logging:} Added colored terminal log output for HMAC tamper detection events (red \texttt{[SECURITY ALERT]} on mismatch) and ECDSA signature verification outcomes (green \texttt{[AUTH VERIFIED]} or red \texttt{[SECURITY ALERT]} per message).

    \item \textbf{Port Decoupling and Multi-Tab Session Sharing:} Eliminated all hardcoded backend port numbers from the frontend, replacing them with \texttt{window.BACKEND\_PORT} injected via the \texttt{/config.js} endpoint. Implemented same-window browser tab session sharing using window identity fingerprinting (\texttt{screenX}, \texttt{screenY}, \texttt{outerWidth}, \texttt{outerHeight}) stored in \texttt{localStorage}. Fixed page-refresh state loss.
\end{itemize}

%------------------------------------------------------------------
\subsection{Sneha Nagmoti}
\textit{GitHub: \texttt{snehanagmoti}}

\textbf{Primary Areas:} Backend Infrastructure, TLS/SSL, Advanced WebSocket Features, Database Attachment Layer, Project Foundation.

\begin{itemize}[leftmargin=*]
    \item \textbf{Project Foundation and Initial Server:} Created the original project scaffold including the FastAPI WebSocket server, \texttt{ConnectionManager}, in-memory broadcast, and the initial HTML/CSS/JS frontend that formed the Lab 3 baseline for the entire team.

    \item \textbf{TLS/SSL Enforcement:} Configured both the backend (\texttt{PORT}=5000) and frontend (\texttt{FRONTEND\_PORT}=3000) Uvicorn servers to use TLS via \texttt{ssl\_keyfile} and \texttt{ssl\_certfile}. Upgraded all connections from \texttt{ws://}/\texttt{http://} to \texttt{wss://}/\texttt{https://}. Created the \texttt{generate\_certs.py} helper script and documented the certificate generation process. Fixed SSL path resolution so the server can be started from any working directory.

    \item \textbf{Certificate Authorization Overlay:} Implemented the \textbf{CONNECTION BLOCKED} in-app overlay that guides users through accepting the self-signed TLS certificate. The overlay links directly to the backend \texttt{/health} URL (added in the same set of changes) and polls until the WebSocket connection resumes automatically.

    \item \textbf{Advanced WebSocket Messaging:} Implemented four new server-side WebSocket message types in a single feature push: soft-delete (\texttt{delete\_message}), whisper routing via \texttt{target\_user}, threaded replies via \texttt{reply\_to}, and \texttt{@mention} detection. Also fixed CORS middleware configuration and cross-platform log encoding issues.

    \item \textbf{Database Attachment Layer and Migrations:} Added the \texttt{attachment} column to the \texttt{messages} table and implemented the \texttt{\_apply\_migrations()} function using \texttt{PRAGMA table\_info} to upgrade the live database schema without data loss. Extended \texttt{save\_message()} to persist file attachment metadata.
    \item \textbf{Complete History Retrieval:} Removed the legacy 50-message constraint from the SQLite fetch queries, allowing the system to serve complete, untruncated conversation histories whenever a user authenticates and joins a chat room.
\end{itemize}

%------------------------------------------------------------------
\subsection{Snehal Suhane}
\textit{GitHub: \texttt{snehalsuhane}}
\begin{itemize}[leftmargin=*]
\item\textbf{Primary Areas:} Multi-Room Lobby, XP Persistence \& Gamification UI, Message Receipts, Interactive Feature Guide, and Code Refactoring.

\item \textbf{Multi-Room Architecture and API Integration:} Architected the backend infrastructure for the multi-room system. Designed the \texttt{rooms} SQLite table to reliably track room metadata, 6-character unique alphanumeric IDs, and creator ownership. Developed the \texttt{GET /rooms}, \texttt{POST /rooms}, and \texttt{GET /rooms/\{id\}} REST endpoints to handle room instantiation and dynamic state fetching for the client.

    \item \textbf{Lobby Dashboard and Room Access UI:} Built a responsive frontend lobby dashboard that dynamically renders active public rooms as interactive cards (displaying custom avatars, names, creators, and real-time live participant counts). Engineered a room creation modal allowing users to customize room aesthetics and toggle public/private visibility. Implemented a dual-access flow where public rooms are joined seamlessly, while private rooms enforce a secure 6-character invite code entry. Finally, added an in-chat room badge to the sidebar to facilitate easy code-sharing among participants.

    \item \textbf{Persistent XP, Ranks, and Badges:} Migrated XP tracking from client-side session state to the server-side \texttt{users} table. Implemented \texttt{db.add\_xp()} for atomic server-side increments. The login response now includes the user's stored XP total so the XP bar is immediately correct without extra requests. Extended the lobby screen to display the user's XP and rank. Added themed rank icons and earned visual badges to both the lobby and chat sidebar. Removed the legacy in-memory statistics object that caused divergence after reconnections.

    \item \textbf{Interactive Feature Guide Dashboard:} Designed and implemented the in-app interactive feature guide. Built a responsive 5-tile grid dashboard modal that visually showcases XP progression, ranks, badges, chat mechanics, and documents the Lab 4 security features.

    \item \textbf{Message Delivery Receipts:} Implemented per-message delivery receipts. The server computes how many room members received each broadcast and returns a \texttt{receipt} event with status \texttt{sent}, \texttt{partial}, or \texttt{delivered\_all}. The sender's message bubble tick indicator updates accordingly.

    \item \textbf{Code Quality, Refactoring, and Stabilization:} Removed dead code, unused variables, and stale debug statements from \texttt{app.js}. Standardised async/await usage and naming conventions. Updated private message usage instructions. Maintained \texttt{.gitignore} hygiene (excluded uploads, certificates, and the database). Resolved multiple Git merge conflicts across \texttt{server.py}, \texttt{db.py}, \texttt{app.js}, and \texttt{index.html} from parallel development branches to stabilize the final build.
\end{itemize}

\section{References}
\subsection*{Acknowledgements}
While AI tools were used to assist with basic document formatting and UI scaffolding, the entirety of the project's architecture, security logic, and backend infrastructure was independently engineered and executed by the team.
\end{document}
